\section{Averaging within a cell}
\label{sec:local}

We next construct restrictions of a cut to one cell.  The restrictions
prescribe its values at a fixed number of terminals and make its values
equal across every internal matching edge.  Averaging a function over
these restrictions has a uniform marginal and, as the cell grows, small
variance over the choice of terminals and matching.  The variance bound
will hold for matrix-valued functions with a constant independent of their
dimensions.

Fix an even integer $d\ge 2$, and let $k$ be a sufficiently large even
integer.  A cell has vertex set $[k]$.  For $b\in\{0,1\}$, let $w$ be the
unique integer in $\{k/2,k/2+1\}$ congruent to $b$ modulo two, and set
\[
  X=\binom{[k]}{w},
  \qquad
  Y_b=\left\{y\in\{0,1\}^d:
       \sum_{\ell=1}^d y_\ell\equiv b\pmod 2\right\},
  \qquad L=|Y_b|=2^{d-1}.
\]
We identify a subset $a\in X$ with its indicator function when referring
to its bits.

Local column data $m$ consist of an ordered list of distinct terminals
$(v_1,\ldots,v_d)$ and a perfect matching of the remaining $k-d$ vertices.
Let $\mathcal M$ be the finite set of these data, with its uniform law.
For $y\in Y_b$, define $\pi_y^m$ to be the uniform probability measure on
the subsets $a\in X$ satisfying
\begin{equation}
  \mathbf 1_a(v_\ell)=y_\ell\quad(1\le\ell\le d),
  \qquad
  \mathbf 1_a(u)=\mathbf 1_a(v)
       \quad\text{for every internal pair }\{u,v\}\text{ of }m.
  \label{local:restrictions}
\end{equation}
To see that this measure is defined, write
\begin{equation}
  s=|y|,\qquad h=\frac{k-d}{2},\qquad
  g=\frac{w-s}{2}.
  \label{local:parameters}
\end{equation}
The parity condition makes $g$ integral.  A subset satisfying
\eqref{local:restrictions} consists of the prescribed terminals and a
choice of $g$ whole pairs among the $h$ internal pairs.  If
$\delta=w-k/2\in\{0,1\}$, then
\begin{equation}
  g=\frac h2+\frac d4+\frac\delta2-\frac s2,
  \qquad
  \left|g-\frac h2\right|\le\frac d4+\frac12.
  \label{local:central}
\end{equation}
Thus $0\le g\le h$ for all sufficiently large $k$, uniformly in $b$ and
$y$.  The support of $\pi_y^m$ has cardinality $\binom hg$.

For a real function $B$ on $X$, define
\[
  (T_yB)(m)=\sum_{a\in X}\pi_y^m(a)B(a).
\]
The same definition applies entrywise to a real matrix-valued function.
Expectations over $a$ and $m$ below use the uniform laws on $X$ and
$\mathcal M$, respectively.

\begin{lemma}[Local averaging]
\label{lem:local-averaging}
For each fixed even $d\ge2$, there are $k_0(d)$ and numbers
$\theta_d(k)\ge0$, defined for even $k\ge k_0(d)$, with
$\theta_d(k)\to0$ as $k\to\infty$, such that the following holds for the
construction above, for both $b\in\{0,1\}$ and every $y\in Y_b$.
First,
\begin{equation}
  \E_m\pi_y^m(a)=\frac1{|X|}\qquad(a\in X).
  \label{local:marginal}
\end{equation}
Second, every real function $B:X\to\R$ satisfies
\begin{equation}
  \E_m\left|(T_yB)(m)-\E_aB(a)\right|^2
  \le\theta_d(k)\,
      \E_a\left|B(a)-\E_aB(a)\right|^2.
  \label{local:scalar-variance}
\end{equation}
For every pair of positive integers $p,q$, the same constant gives, for
every $B:X\to\R^{p\times q}$,
\begin{equation}
  \E_m\normF{(T_yB)(m)-\E_aB(a)}^2
  \le\theta_d(k)\,
      \E_a\normF{B(a)-\E_aB(a)}^2
  \le\theta_d(k)\,\E_a\normF{B(a)}^2.
  \label{local:matrix-variance}
\end{equation}
One may take $\theta_d(k)=O_d(k^{-1/4}\sqrt{\log(k+2)})$.
\end{lemma}

\begin{proof}
The uniform law of $m$ and the restrictions
\eqref{local:restrictions} are equivariant under permutations of the
$k$ vertices.  Since the symmetric group is transitive on $X$, the
marginal law of $a$, obtained by first sampling $m$ and then
$a\sim\pi_y^m$, is uniform.  This proves \eqref{local:marginal}.

Fix $b$ and $y$, and write $T=T_y$.  Equip the function spaces on $X$
and $\mathcal M$ with the inner products for uniform probability, and
put $K=T^*T$.  This is a positive semidefinite self-adjoint operator.
The marginal identity gives $K1=1$.  We will show that
$\tr(K^2)=o(k)$, while every nonconstant eigenspace has dimension at
least $k/2$.  Together these estimates force the largest eigenvalue
on the orthogonal complement of the constants to tend to zero.

\emph{The two-sample kernel.}
Sample $m$ uniformly, and then sample $a,a'$ independently from
$\pi_y^m$ conditional on $m$.  Let $J(a,a')$ be their joint
probability, and let $N=|X|$.  The density of this joint law relative
to two independent uniform elements of $X$ is
$D(a,a')=N^2J(a,a')$.  Indeed, for scalar functions $B_1,B_2$,
\[
  \langle B_1,KB_2\rangle
   =\sum_{a,a'\in X}J(a,a')B_1(a)B_2(a').
\]
Consequently the ordinary matrix entries of $K$ are $D(a,a')/N$, and
symmetry gives
\begin{equation}
  \tr(K^2)=\frac1{N^2}\sum_{a,a'\in X}D(a,a')^2.
  \label{local:density-trace}
\end{equation}

The joint law is permutation-invariant, so $D(a,a')$ depends only on
$|a\cap a'|$.  Conditional on $m$, the two subsets choose independent
$g$-subsets of the $h$ internal pairs.  If the second choice replaces
exactly $j$ pairs of the first, their intersection has size
$s+2(g-j)=w-2j$.  The probability of this event is
\[
  P_j=\frac{\binom gj\binom{h-g}j}{\binom hg},
  \qquad 0\le j\le g_*:=\min(g,h-g).
\]
For two independent uniform elements of $X$, the probability of the
same intersection size is
\[
  Q_j=\frac{\binom w{2j}\binom{k-w}{2j}}{\binom kw}.
\]
Both laws are uniform within each intersection class.  The density
there is therefore $P_j/Q_j$, and it vanishes on the remaining
classes.  The identities
$w=2g+s$ and $k-w=2(h-g)+(d-s)$ ensure that every displayed $Q_j$
is positive.  Equation~\eqref{local:density-trace} becomes
\begin{equation}
  \tr(K^2)
   =\sum_{j=0}^{g_*}\frac{P_j^2}{Q_j}
   =\frac{\binom kw}{\binom hg^2}
      \sum_{j=0}^{g_*}
      \frac{\binom gj^2\binom{h-g}j^2}
           {\binom w{2j}\binom{k-w}{2j}}.
  \label{local:trace-formula}
\end{equation}

\emph{A trace estimate.}
Since $k=2h+d$, $w=k/2+O(1)$, and
$g=h/2+O_d(1)$, the central binomial estimates give
\begin{equation}
  \frac{\binom kw}{\binom hg^2}=O_d(\sqrt h),
  \label{local:central-ratio}
\end{equation}
uniformly in $b$ and $y$.  Moreover,
$\binom w{2j}\ge\binom{2g}{2j}$ and
$\binom{k-w}{2j}\ge\binom{2(h-g)}{2j}$.
For integers $0\le j\le u$, the identity
\[
  \frac{\binom uj^2}{\binom{2u}{2j}}
  =\frac{\binom{2j}j\binom{2(u-j)}{u-j}}{\binom{2u}u}
\]
and the bounds $\binom{2v}v\asymp4^v/\sqrt{v+1}$ imply
\begin{equation}
  \frac{\binom uj^2}{\binom{2u}{2j}}
   \le C\sqrt{\frac{u+1}{(j+1)(u-j+1)}}
  \label{local:binomial-ratio}
\end{equation}
with an absolute constant $C$, including at the endpoints $j=0,u$.
Applying this bound with $u=g$ and $u=h-g$ shows that the sum in
\eqref{local:trace-formula} is at most
\[
  C^2\sum_{j=0}^{g_*}
    \frac{h+1}{(j+1)(g_*-j+1)}
  =\frac{2C^2(h+1)}{g_*+2}\sum_{j=1}^{g_*+1}\frac1j.
\]
By \eqref{local:central}, $g_*=h/2+O_d(1)$, so this is
$O_d(\log(h+2))$.  Thus
\begin{equation}
  \tr(K^2)=O_d(\sqrt h\log(h+2)).
  \label{local:trace-bound}
\end{equation}

\emph{From trace to contraction.}
We now use permutation symmetry to ensure that every nonconstant
eigenvalue contributes many times to the trace.  This is the
multiplicity principle used, in a different setting, in
\cite[Lemma~3.2]{gowers}; here an elementary argument suffices.

We claim that, for even $k\ge6$, every nonzero permutation-invariant
subspace $V\subseteq\R^X$ orthogonal to the constants has dimension
at least $k/2$.  Choose $k/2$ disjoint transpositions of the vertices.
Their actions on $V$ are commuting orthogonal involutions, so they
admit simultaneous diagonalization with sign patterns in
$\{1,-1\}^{k/2}$.  If a nonzero simultaneous eigenspace has a pattern
with $j$ negative signs, where $0<j<k/2$, permutations of the vertex
pairs give $\binom{k/2}j$ distinct patterns.  Their eigenspaces are
linearly independent, whence
$\dim V\ge\binom{k/2}j\ge k/2$.

Otherwise every occurring pattern is constant, and the product of
any two of the chosen transpositions acts as the identity on $V$.
Invariance of $V$ under all permutations implies the same for every
product of two disjoint transpositions.  Given five distinct labels
$a,b,c,e,f$, the identity
\[
  [(ab)(ef)]\,[(bc)(ef)]=(ab)(bc)
\]
then shows that every three-cycle acts as the identity.  Thus the
alternating group fixes $V$ pointwise.  It is transitive on $X$:
a permutation taking one $w$-subset to another can have its parity
corrected by a transposition inside the latter subset or its
complement.  Such a transposition exists for the present values of
$k$ and $w$.  Every function on $X$ invariant under the alternating
group is therefore constant, contrary to the choice of $V$.
This proves the claim.

The permutation invariance of the kernel means that $K$ commutes
with every vertex permutation.  Each full eigenspace of its
restriction to the orthogonal complement of the constants is
therefore an invariant subspace of the kind just considered.
Writing $\mu$ for the largest eigenvalue on that complement,
positive semidefiniteness and
\eqref{local:trace-bound} imply
\[
  \frac k2\mu^2\le\tr(K^2)-1,
  \qquad
  \mu=O_d\bigl(k^{-1/4}\sqrt{\log(k+2)}\bigr).
\]
For $B_0=B-\E_aB(a)$, it follows that
\[
  \E_m|(TB)(m)-\E_aB(a)|^2
   =\langle B_0,KB_0\rangle
   \le\mu\,\E_a|B_0(a)|^2.
\]
All estimates were uniform in $b$ and $y$, so a single choice of
$\theta_d(k)$ proves \eqref{local:scalar-variance}.  Summing that
inequality over the entries of a matrix proves
\eqref{local:matrix-variance}, with exactly the same constant.
\end{proof}

The matrix estimate also applies after fixing auxiliary data, provided
that the conditional law of $m$ given those data is uniform on
$\mathcal M$.  The data may specify any matrix-valued function $B$
on $X$: fixing them must determine the whole function before $m$ is
sampled, but the dependence of $B(a)$ on $a$ is unrestricted.
The lemma then applies separately at each fixing.  Thus conditioning
on other cells may select an arbitrary function on the present cell;
averaging the resulting inequalities introduces no factor depending
on its matrix dimensions.
